% ==================================================================
% id:       comb-006
% topic:    קומבינטוריקה
% subtopic: בחירה עם חזרות
% solution: full
% ==================================================================
בכיתה 20 תלמידים.
\begin{enumerate}
\item בכמה דרכים אפשר לבחור ועד בן 3 חברים?
\begin{solbox}
$\binom{20}3=\ans{1140}$.
\end{solbox}
\item בכמה דרכים אפשר לבחור יו"ר, סגן וגזבר (שלושה אנשים שונים)?
\begin{solbox}
כאן יש חשיבות לסדר (מי יו"ר, מי סגן ומי גזבר), ואין חזרות (אותו אדם לא ממלא שני תפקידים). בוחרים יו"ר מתוך 20 תלמידים, אחר כך סגן מתוך 19 הנותרים, ואחר כך גזבר מתוך 18 הנותרים. לפי עקרון הכפל יש
\[
20\cdot19\cdot18=\frac{20!}{17!}=\ans{6840}
\]
דרכים.

\textbf{דרך נוספת:} קודם בוחרים את שלושת חברי הוועד, בלי תפקידים: $\binom{20}3=1140$ דרכים (סעיף א). אחר כך מחלקים ביניהם את שלושת התפקידים: $3!=6$ דרכים. לפי עקרון הכפל, $1140\cdot6=6840$.
\end{solbox}
\item במשך שלושה ימים ברצף, בכל יום בוחרים תלמיד תורן. אותו תלמיד יכול להיות תורן ביותר מיום אחד. כמה שיבוצים יש?
\begin{solbox}
בכל אחד משלושת הימים יש 20 אפשרויות, בלי תלות במה שנבחר בימים האחרים. כאן יש חשיבות לסדר (מי תורן באיזה יום) ויש חזרות, ולכן לפי עקרון הכפל יש $20^3=\ans{8000}$ שיבוצים.
\end{solbox}
\item הסבירו במילים למה התשובה ל-(ב) היא בדיוק $3!$ כפול התשובה ל-(א).
\begin{solbox}
כל ועד אפשר לסדר בתפקידים ב-$3!$ דרכים.
\end{solbox}
\end{enumerate}
